% Generated deterministically from the audited Markdown source.
% Responsible author: Carptopus.
\documentclass[11pt]{article}
\usepackage[a4paper,margin=25mm]{geometry}
\usepackage[T1]{fontenc}
\usepackage[utf8]{inputenc}
\usepackage{lmodern}
\usepackage{microtype}
\usepackage{amsmath,amssymb,amsthm,mathtools}
\usepackage{needspace}
\usepackage{xcolor}
\usepackage[colorlinks=true,linkcolor=blue!55!black,citecolor=blue!55!black,urlcolor=blue!65!black]{hyperref}
\usepackage{xurl}
\usepackage{fancyhdr}

\newtheorem{theorem}{Theorem}
\newtheorem{lemma}[theorem]{Lemma}

\pagestyle{fancy}
\fancyhf{}
\fancyhead[L]{Carptopus}
\fancyhead[R]{Rank-two matroid $h^*$-polynomials}
\fancyfoot[C]{\thepage}
\setlength{\headheight}{14pt}
\setlength{\parindent}{0pt}
\setlength{\parskip}{0.42em}
\setlength{\emergencystretch}{4em}
\urlstyle{same}

\title{Real-rooted $h^*$-polynomials of rank-two matroids\\
with parallel classes of size at most three}
\author{Carptopus\\\href{mailto:carptopus@163.com}{carptopus@163.com}}
\date{4 September 2026}

\hypersetup{
  pdftitle={Real-rooted h-star polynomials of rank-two matroids with parallel classes of size at most three},
  pdfauthor={Carptopus},
  pdfsubject={Ehrhart h-star real-rootedness for rank-two matroid base polytopes},
  pdfkeywords={matroid base polytope, Ehrhart h-star polynomial, real-rootedness, rank-two matroid, parallel class}
}

\begin{document}
\maketitle

\begin{center}
\fcolorbox{black!35}{black!4}{\parbox{0.88\textwidth}{\centering\small
Public Beta manuscript, version 0.1-beta. The mathematical proof and
complete manuscript have passed project-internal zero-trust audits. External
mathematical review and formal peer review are pending.}}
\end{center}

\begin{abstract}


Ferroni, Jochemko and Schröter proved that the Ehrhart $h^*$-polynomial of
the base polytope of a rank-two sparse paving matroid is real-rooted. After
loops are deleted, their rank-two hypothesis says that every parallel class
has size at most two. We extend their result to parallel classes of size at
most three. The
new ingredient is a uniform estimate for the cubic correction term
$p^*_{3,n}$ in their explicit $h^*$-formula. After the substitution
$x=-\tan^2\theta$, this estimate preserves the alternating signs used in
their proof. We prove the estimate analytically for $n\ge 10$, settle
$n=9$ by an exact Sturm certificate, and check the finitely many smaller
connected and disconnected cases exactly. Consequently the $h^*$-vector is
log-concave and unimodal. We also identify the necessary Laguerre-polynomial
inequality suggested by the scaling limit of the same method for an
arbitrary fixed bound on parallel-class size.

\end{abstract}

\section{Introduction}


For a matroid $M$ on a finite ground set, let $P(M)$ denote its base
polytope and let $h^*(P(M),x)$ be the numerator of its Ehrhart series. Ferroni
conjectured that the $h^*$-polynomial of every matroid base polytope is
real-rooted \cite{ferroni-minimal}. Ferroni, Jochemko and Schröter subsequently proved this for
rank-two sparse paving matroids \cite{fjs}. They also gave an explicit formula for
every rank-two matroid and proposed a precise analytic estimate as a possible
route beyond the sparse paving case.

In a loopless rank-two matroid the rank-one flats are exactly the parallel
classes. Sparse paving is therefore equivalent in this setting to requiring
every parallel class to have size at most two. Our main result advances this
boundary by one full structural layer.

\Needspace{0.16\textheight}
\begin{theorem}
\label{thm:main}

Let $M$ be a rank-two matroid. If every nonloop parallel class of $M$ has
size at most three, then $h^*(P(M),x)$ has only real zeros. Consequently its
coefficient sequence is log-concave and unimodal.
\end{theorem}

The proof is computer-assisted only at the finite boundary $n\le 9$: the
$n=9$ estimate is certified by exact Sturm arithmetic, and the 31 connected
types on at most eight nonloop elements and the six disconnected types are
checked exactly. All larger ground sets are covered by one analytic argument.

The theorem does not settle the conjecture for all rank-two matroids. The
largest parallel class is bounded by three, whereas a general rank-two matroid
may have a parallel class of arbitrary size. Section 6 records the exact
uniform inequality that remains for extending the method.

\section{The rank-two formula}


Deleting the loops of $M$ removes coordinates that are identically zero on
$P(M)$, so it preserves the lattice-isomorphism type of the base polytope.
We may therefore assume that $M$ is loopless and write $n$ for the number
of its remaining elements.

Suppose first that $M$ is connected. Its parallel classes form a partition
of the ground set, say with sizes $a_1,\ldots,a_s$, where $s\ge 3$ and

\[
a_1+\cdots+a_s=n.
\]

Ferroni, Jochemko and Schröter \cite[Proposition~5.1]{fjs} proved

\[
h^*(P(M),x)
=\sum_{j=0}^{\lfloor n/2\rfloor}\binom{n}{2j}x^j
-\sum_{i=1}^s p^*_{a_i,n}(x),                                      \tag{2.1}
\]

where

\[
p^*_{a,n}(x)
=\sum_{k=1}^{a}\sum_{j\ge1}
  \binom{k}{j}\binom{n-k-1}{j-1}x^j.                               \tag{2.2}
\]

They established \cite[Lemma~5.3]{fjs}, for $a\in\{1,2\}$, $n\ge9$, and
$0<y\le1$,

\[
\left|y^{n/2}p^*_{a,n}\!\left(\frac{y-1}{y}\right)\right|<\frac an.
                                                                        \tag{2.3}
\]

Our new ingredient is the next case.

\Needspace{0.16\textheight}
\begin{lemma}
\label{lem:a3}

For every $n\ge9$ and $0<y\le1$,
\end{lemma}

\[
\left|y^{n/2}p^*_{3,n}\!\left(\frac{y-1}{y}\right)\right|<\frac3n.
                                                                        \tag{2.4}
\]

Directly from (2.2),

\[
p^*_{3,n}(x)
=6x+(4n-15)x^2+\frac{(n-4)(n-5)}2x^3.                                \tag{2.5}
\]

We prove Lemma 2.1 in the next two sections.

\section{The uniform estimate for \texorpdfstring{$n\ge10$}{n >= 10}}


Put

\[
u=n\frac{1-y}{y}\ge0,
\qquad y=\left(1+\frac un\right)^{-1}.
\]

After substituting (2.5), inequality (2.4) is equivalent to

\[
|P_n(u)|<3\left(1+\frac un\right)^{n/2},                              \tag{3.1}
\]

where

\[
P_n(u)=-6u+\left(4-\frac{15}{n}\right)u^2
-\frac12\left(1-\frac9n+\frac{20}{n^2}\right)u^3.                  \tag{3.2}
\]

Let $T_{5,n}(u)$ be the Taylor polynomial of degree five at zero of
$(1+u/n)^{n/2}$. For $n\ge10$, its sixth-order Taylor remainder is
nonnegative on $u\ge0$; this includes $n=10$, when the remainder vanishes.
Thus

\[
\left(1+\frac un\right)^{n/2}\ge T_{5,n}(u).                         \tag{3.3}
\]

It is therefore enough to prove
$3T_{5,n}(u)-P_n(u)>0$ and $3T_{5,n}(u)+P_n(u)>0$.

For the upper inequality, direct expansion gives

\[
3T_{5,n}(u)-P_n(u)
=3+\frac{15}{2}u-A_nu^2+B_nu^3+E_nu^4+F_nu^5,                       \tag{3.4}
\]

where

\[
A_n=\frac{29n-114}{8n},\qquad
B_n=\frac{3(n-4)(3n-14)}{16n^2},
\]

\[
E_n=\frac{(n-6)(n-4)(n-2)}{128n^3},\qquad
F_n=\frac{(n-8)(n-6)(n-4)(n-2)}{1280n^4}.
\]

All four quantities are positive for $n\ge10$, and

\[
30B_n-A_n^2
=\frac{239n^2-2748n+7164}{64n^2}>0.                              \tag{3.5}
\]

Indeed, on writing $n=9+v$, the numerator in (3.5) becomes
$239v^2+1554v+1791$. For $u>0$, the arithmetic--geometric mean
inequality and (3.5) yield

\[
\frac{15}{2}u+B_nu^3
\ge 2\sqrt{\frac{15B_n}{2}}u^2>A_nu^2.
\]

Equation (3.4) is therefore positive; at $u=0$ this is immediate.

For the lower inequality, write

\[
3T_{5,n}(u)+P_n(u)=3-\frac92u+u^2R_n(u),                             \tag{3.6}
\]

where

\[
R_n(u)=C_n-D_nu+E_nu^2+F_nu^3,
\]

\[
C_n=\frac{7(5n-18)}{8n},\qquad
D_n=\frac{(n-4)(7n-38)}{16n^2}.                                    \tag{3.7}
\]

The discriminant of the cubic $R_n-1$ is

\[
-\frac{(n-6)(n-4)^2(n-2)}{104857600n^{10}}Q(n),                    \tag{3.8}
\]

and, with $n=9+v$,

\[
\begin{aligned}
Q(9+v)={}&8703v^6+204112v^5+1743391v^4+6933784v^3\\
&+13118701v^2+10391512v+2399413.
\end{aligned}
\]

Hence the discriminant in (3.8) is negative. The cubic $R_n-1$ has
positive leading coefficient and

\[
(R_n-1)(0)=\frac{27n-126}{8n}>0.
\]

It has exactly one real zero, which must therefore be negative. Consequently

\[
R_n(u)>1\qquad (u\ge0).                                             \tag{3.9}
\]

If $u\ge37/10$, equations (3.6) and (3.9) give

\[
3T_{5,n}(u)+P_n(u)>3-\frac92u+u^2>0,
\]

because the final quadratic is increasing on this interval and equals
$1/25$ at $u=37/10$.

If $0\le u\le37/10$, then $C_n\ge21/8$, $D_n\le7/16$, and
$E_n,F_n\ge0$. Thus

\[
3T_{5,n}(u)+P_n(u)
\ge q(u):=3-\frac92u+\frac{21}{8}u^2-\frac7{16}u^3.                 \tag{3.10}
\]

The two stationary points of $q$ are $2\pm2\sqrt7/7$, where its
values are $1\mp\sqrt7/7$. Together with
$q(0)=3$ and $q(37/10)=2009/16000$, this proves that $q$ is positive
on the entire interval. Equations (3.3)--(3.10) prove (3.1) for $n\ge10$.

\section{The finite boundary}


For $n=9$, let $D=(1+u/9)^{9/2}$. Exact expansion gives

\[
9D^2-P_9(u)^2=\frac{N(u)}{43046721},                                \tag{4.1}
\]

where

\[
\begin{aligned}
N(u)={}&u^9+81u^8+2916u^7-594864u^6+25627266u^5\\
&-290698227u^4+1249949232u^3-1377495072u^2\\
&+387420489u+387420489.
\end{aligned}
\]

The exact Sturm chain of $N$ has degree profile

\[
(9,8,6,5,4,3,2,1,0).
\]

Its signs at zero and at positive infinity are respectively

\[
(+,+,0,-,-,+,-,-,-),
\qquad (+,+,+,-,-,-,+,+,-).
\]

After zeros are omitted, both sequences have three sign variations. Hence
$N$ has no positive zero. Since $N(0)>0$, equation (4.1) is positive for
all $u\ge0$, proving $|P_9(u)|<3D$ and therefore Lemma 2.1 for $n=9$.

For $3\le n\le8$, a connected loopless rank-two matroid is determined up
to isomorphism by an integer partition of $n$ into at least three parts.
Restricting each part to size at most three leaves, for
$n=3,4,5,6,7,8$, respectively

\[
1,\ 2,\ 4,\ 6,\ 8,\ 10
\]

possibilities. The accompanying exact verifier constructs (2.1) for all 31
partitions, factors off repeated factors, and uses Sturm root counts to place
every zero, with multiplicity, in $(-\infty,0)$. It also uses SymPy's exact
algebraic-root routine as a second check within the same computer-algebra
system.

\section{Proof of the main theorem}

\begin{proof}[Proof of Theorem~\ref{thm:main}]

Assume first that $M$ is connected and $n\ge9$. Set

\[
x=-\tan^2\theta,\qquad 0\le\theta<\frac{\pi}{2}.
\]

The classical identity

\[
\sum_j\binom n{2j}(-\tan^2\theta)^j
=\frac{\cos(n\theta)}{\cos^n\theta}                                \tag{5.1}
\]

turns (2.1) into

\[
h^*(P(M),-\tan^2\theta)=\frac{q(\theta)}{\cos^n\theta},             \tag{5.2}
\]

where

\[
q(\theta)=\cos(n\theta)
-\sum_{i=1}^s\cos^n\theta\,
  p^*_{a_i,n}(-\tan^2\theta).                                      \tag{5.3}
\]

Put $y=\cos^2\theta$. For $a_i=1,2$, use (2.3), and for $a_i=3$,
use Lemma 2.1. Since the $a_i$ sum to $n$, equations (5.3) and the strict
estimates give

\[
|q(\theta)-\cos(n\theta)|
<\sum_{i=1}^s\frac{a_i}{n}=1.                                     \tag{5.4}
\]

At the points $\theta_j=j\pi/n$, equation (5.4) makes $q(\theta_j)$
strictly positive for even $j$ and strictly negative for odd $j$.

When $n$ is odd, all points with
$0\le j\le\lfloor n/2\rfloor$ lie below $\pi/2$. When $n=2m$ is
even, the last point is the endpoint $\pi/2$, outside the domain of (5.4).
In this case every correction in (5.3) is a linear combination of terms

\[
(-1)^j\sin^{2j}\theta\cos^{n-2j}\theta,
\qquad 1\le j\le3.
\]

Since $n\ge10$, the exponent $n-2j$ is at least four. All corrections
therefore tend to zero at the endpoint, and $q$ extends continuously with

\[
q(\pi/2)=\cos(n\pi/2)=(-1)^m.
\]

Thus, in both parity cases, $q$ has a zero in each of the
$\lfloor n/2\rfloor$ consecutive open intervals between the sign-test
points. Under the injective map
$\theta\mapsto-\tan^2\theta$, these become distinct negative zeros of
$h^*(P(M),x)$.

The first sum in (2.1) has degree $\lfloor n/2\rfloor$, while every
correction has degree at most three. For $n\ge9$, the former degree is at
least four, so its positive leading term cannot be cancelled. Hence the
number of zeros found equals the degree, and all zeros are real. The smaller
connected cases were settled in Section 4.

It remains to treat disconnected matroids. After loops are deleted, a
disconnected rank-two matroid is $U_{1,a}\oplus U_{1,b}$, where $a,b\ge1$.
Under our hypothesis $a,b\le3$, and $P(M)$ is the product
$\Delta_{a-1}\times\Delta_{b-1}$. For the six unordered pairs

\[
(1,1),(1,2),(1,3),(2,2),(2,3),(3,3),
\]

the corresponding $h^*$-polynomials are

\[
1,\quad1,\quad1,\quad1+x,\quad1+2x,\quad1+4x+x^2.
\]

They are all real-rooted. This completes the proof of Theorem 1.1. The
log-concavity and unimodality conclusions follow from Newton's inequalities
and the nonnegativity of Ehrhart $h^*$-coefficients.
\end{proof}


\section{The boundary of the method}


Ferroni, Jochemko and Schröter \cite[Section~6.1]{fjs} suggested proving, for every
fixed $a$, that for all sufficiently large $n$,

\[
\left|y^{n/2}p^*_{a,n}\!\left(\frac{y-1}{y}\right)\right|<\frac an.
                                                                        \tag{6.1}
\]

The same sign-interlacing argument would then cover rank-two matroids whose
parallel classes have size at most $a$, subject to the degree condition.
A useful necessary asymptotic test is visible in the scaling used above. For
fixed $a$, the hockey-stick identity gives

\[
\lim_{n\to\infty} n\,p^*_{a,n}(-u/n)
=\sum_{j=1}^a(-1)^j\binom{a+1}{j+1}\frac{u^j}{(j-1)!}
=-uL_{a-1}^{(2)}(u),                                                 \tag{6.2}
\]

where $L_m^{(2)}$ is a generalized Laguerre polynomial. Since
$(1+u/n)^{-n/2}\to e^{-u/2}$, pointwise passage to the limit in (6.1)
would require the non-strict inequality

\[
\sup_{u\ge0}|u e^{-u/2}L_{a-1}^{(2)}(u)|\le a.                       \tag{6.3}
\]

Numerical evidence suggests the stronger strict inequality, which would give
a useful pointwise margin. Even that strict bound would not by itself imply
the finite-$n$ estimate (6.1) on the unbounded interval $u\ge0$: a uniform
convergence argument and separate tail control would still be needed.
Standard Laguerre bounds are too weak after multiplication by $u$, and
treating $a=4,5,\ldots$ by separate increasingly high-degree polynomial
certificates would not supply a uniform mechanism. The present note therefore
stops at the natural $a=3$ boundary.

\section{Verification and reproducibility}


The supplementary Python script uses integer and rational arithmetic through
SymPy. It performs the following checks:

\begin{enumerate}
\item derives (2.5) directly from (2.2);
\item reconstructs both Taylor expansions in Section 3 and verifies their
coefficient, discriminant, stationary-point and endpoint certificates;
\item reconstructs $N(u)$ in (4.1), its Sturm profile and endpoint signs;
\item includes a positive-root negative control for the Sturm routine;
\item checks all 31 small connected partitions using exact Sturm counts and
an additional exact algebraic-root check in the same SymPy stack;
\item derives and checks all six disconnected cases.
\end{enumerate}

All failure conditions are explicit exceptions, not Python \texttt{assert}
statements, so the checks remain active under \texttt{python -O}. The script is a
reproducibility aid and supplies the exact finite portion of the proof; it
does not replace the analytic argument for $n\ge10$.

The audited environment used Python 3.13.11 and SymPy 1.14.0. The SHA-256 of
the audited script is
\texttt{00FA0B2A2A7481B6B7E562EDC5EFF2DD4DD63ED776F3BA1E39FCFCED7EFEF9C9}.

\section{Prior-work and novelty boundary}


Ferroni, Jochemko and Schröter \cite{fjs} supplied both the general rank-two formula
and the sign-interlacing method, and proved the case in which parallel classes
have size at most two. They explicitly identified (6.1) as a possible route
to any fixed upper bound. Our contribution is limited to proving the next
structural case $a=3$, including all finite and disconnected boundaries, and
to identifying the Laguerre scaling limit and the additional uniformity
obligations for a parameter-uniform extension.

Jiang \cite{jiang} later gave combinatorial formulas for the $h^*$-polynomials of all
positroid polytopes. Since every rank-two matroid is isomorphic to a positroid
\cite[Corollary~6.2]{fjs}, that work gives a broader formula-level context, but it
does not state the real-rootedness theorem proved here.

Targeted searches completed on 4 September 2026 found no paper directly
covering parallel classes of size at most three. This supports candidate
novelty but does not establish global priority. We do not claim to settle
real-rootedness for all rank-two matroids or for all matroid base polytopes.

\section*{AI-assisted research disclosure}


Carptopus is the responsible author and takes responsibility for the claims,
proofs, source use and released artifacts. OpenAI Codex was used for
AI-assisted literature organization, proof development and stress testing,
exact verification, adversarial auditing and manuscript preparation.

\begin{thebibliography}{9}
\footnotesize
\raggedright

\bibitem{ferroni-minimal}
L.~Ferroni,
\emph{On the Ehrhart polynomial of minimal matroids},
Discrete \& Computational Geometry \textbf{68} (2022), 255--273.
\href{https://doi.org/10.1007/s00454-021-00313-4}{doi:10.1007/s00454-021-00313-4}.

\bibitem{fjs}
L.~Ferroni, K.~Jochemko and B.~Schr\"oter,
\emph{Ehrhart polynomials of rank two matroids},
Advances in Applied Mathematics \textbf{141} (2022), 102410.
\href{https://doi.org/10.1016/j.aam.2022.102410}{doi:10.1016/j.aam.2022.102410}.

\bibitem{jiang}
Y.~Jiang,
\emph{The Ehrhart $h^*$-polynomials of positroid polytopes},
arXiv:2410.01743v3 (2025).
\url{https://arxiv.org/abs/2410.01743}.

\end{thebibliography}

\end{document}
